\PassOptionsToPackage{table}{xcolor}
\documentclass[10pt,aspectratio=169]{beamer}
\newcommand{\LectureNumber}{07}
\input{course-theme.tex}
\title{Boolean expressions}
\subtitle{CSC103 Programming Fundamentals / Lecture 07}
\author{Dr. Muhammad Farhan}
\institute{Associate Professor of Computer Science\\Department of Computer Science\\COMSATS University Islamabad\\Sahiwal Campus}
\date{Fall 2026}
\begin{document}
\begin{frame}\titlepage\end{frame}

\begin{frame}[fragile,t]{The problem for this lecture}
\begin{columns}[T,onlytextwidth]\begin{column}{0.60\textwidth}
A ticket discount applies to ages below 12 or at least 60. Write the condition and test ages 11, 12, 59 and 60.
\end{column}\begin{column}{0.35\textwidth}\small
Express each qualifying age group as its own comparison.
\end{column}\end{columns}
\note{Introduce the target problem before explaining the tools. Revisit it in the guided solution later.\par Syllabus reading: Deitel Ch 7 (as listed). CLO-2.}
\end{frame}

\begin{frame}[fragile,t]{Learning objectives}
\begin{columns}[T,onlytextwidth]\begin{column}{0.60\textwidth}
\begin{itemize}
\item Construct relational and logical conditions
\item Evaluate Boolean expressions
\item Choose explicit grouping for compound rules
\end{itemize}
\end{column}\begin{column}{0.35\textwidth}\small
CLO-2

Relational operators\par Logical combinations\par Intervals and precedence\par Negation and boundary cases
\end{column}\end{columns}
\note{Explain how each objective supports the opening problem. The final questions check these objectives.\par Syllabus reading: Deitel Ch 7 (as listed). CLO-2.}
\end{frame}

\begin{frame}[fragile,t]{Relational operators}
\begin{itemize}
\item \textless{}, \textless{}=, \textgreater{} and \textgreater{}= compare numeric values.
\item == tests primitive equality and != tests inequality.
\item A relational expression produces true or false.
\end{itemize}
\vspace{1mm}\begin{block}{Example}
\begin{lstlisting}
int marks = 50;
boolean passed = marks >= 50;
boolean perfect = marks == 100;
\end{lstlisting}
\end{block}
\note{For Strings, content comparison uses equals and is taught later. Do not treat assignment = as an equality test.\par Syllabus reading: Deitel Ch 7 (as listed). CLO-2.}
\end{frame}

\begin{frame}[fragile,t]{Logical combinations}
\begin{itemize}
\item \&\& requires both conditions to be true.
\item || requires at least one condition to be true.
\item ! reverses a Boolean value.
\end{itemize}
\vspace{1mm}\begin{block}{Example}
\begin{lstlisting}
boolean eligible = age >= 18 && registered;
boolean weekend = day == 6 || day == 7;
\end{lstlisting}
\end{block}
\note{Use four combinations of two Boolean inputs to construct AND and OR truth tables on the board.\par Syllabus reading: Deitel Ch 7 (as listed). CLO-2.}
\end{frame}

\begin{frame}[fragile,t]{Intervals and precedence}
\begin{itemize}
\item Write a range check as two comparisons joined with \&\&.
\item Relational tests bind before \&\&, which binds before ||.
\item Parentheses make mixed rules easier to inspect.
\end{itemize}
\vspace{1mm}\begin{block}{Example}
\begin{lstlisting}
boolean valid = marks >= 0 && marks <= 100;
boolean accepted = (a || b) && c;
\end{lstlisting}
\end{block}
\note{0 \textless{}= marks \textless{}= 100 is not valid Java. Explain how the first comparison becomes a boolean, which cannot participate in the second numeric comparison.\par Syllabus reading: Deitel Ch 7 (as listed). CLO-2.}
\end{frame}

\begin{frame}[fragile,t]{Negation and boundary cases}
\begin{itemize}
\item The opposite of x \textless{} limit is x \textgreater{}= limit.
\item De Morgan: !(a \&\& b) equals !a || !b.
\item Tests at the boundary reveal missing equality cases.
\end{itemize}
\vspace{1mm}\begin{block}{Example}
\begin{lstlisting}
!(age >= 18 && registered)

Equivalent:
age < 18 || !registered
\end{lstlisting}
\end{block}
\note{Ask for values immediately below, at and above a threshold. Truth tables can verify transformations without relying on intuition.\par Syllabus reading: Deitel Ch 7 (as listed). CLO-2.}
\end{frame}

\begin{frame}[fragile,t]{Translating a sentence into a condition}
\begin{itemize}
\item "Below 12" uses \textless{} 12, while "at least 60" uses \textgreater{}= 60.
\item The word "or" permits either qualifying group.
\item Test the endpoints to confirm that the inequalities reflect the wording.
\end{itemize}
\vspace{1mm}\begin{block}{Example}
\begin{lstlisting}
Discount rule:
age < 12 || age >= 60
\end{lstlisting}
\end{block}
\note{Explain the mechanism using the displayed example. "Below 12" uses \textless{} 12, while "at least 60" uses \textgreater{}= 60. The word "or" permits either qualifying group. Test the endpoints to confirm that the inequalities reflect the wording.\par Syllabus reading: Deitel Ch 7 (as listed). CLO-2.}
\end{frame}

\begin{frame}[fragile,t]{A Boolean expression is a value}
\begin{itemize}
\item The result of a comparison can be stored, passed or combined.
\item A condition must ultimately produce boolean in Java.
\item A number such as 1 is not a replacement for true.
\end{itemize}
\vspace{1mm}\begin{block}{Example}
\begin{lstlisting}
boolean discounted = age < 12 || age >= 60;
System.out.println(discounted);
\end{lstlisting}
\end{block}
\note{Explain the mechanism using the displayed example. The result of a comparison can be stored, passed or combined. A condition must ultimately produce boolean in Java. A number such as 1 is not a replacement for true.\par Syllabus reading: Deitel Ch 7 (as listed). CLO-2.}
\end{frame}


\begin{frame}[fragile,t]{Build a compound condition from two facts}
\begin{center}\begin{tikzpicture}
\node[box] (a) at (0,0) {age \textgreater{}= 18};
\node[box] (b) at (0,-1.6) {hasCard};
\node[decision] (c) at (4,-0.8) {AND};
\node[alt] (d) at (8,-0.8) {admitted};
\draw[arr] (a.east)--++(.6,0)|-(c.west);\draw[arr] (b.east)--++(.6,0)|-(c.west);\draw[arr] (c)--(d);
\end{tikzpicture}\end{center}
\begin{block}{What to notice}Admission requires both facts to be true. One true input is insufficient.\end{block}
\note{Trace each labelled connection. Admission requires both facts to be true. One true input is insufficient.}
\end{frame}
\begin{frame}[fragile,t]{Key distinctions}
\small\renewcommand{\arraystretch}{1.5}
\rowcolors{2}{pale}{white}
\begin{tabularx}{\textwidth}{>{\raggedright\arraybackslash}X>{\raggedright\arraybackslash}X>{\raggedright\arraybackslash}X>{\raggedright\arraybackslash}X}
\rowcolor{navy}
\color{white}\bfseries A & \color{white}\bfseries B & \color{white}\bfseries A \&\& B & \color{white}\bfseries A || B\\
false & false & false & false\\
false & true & false & true\\
true & false & false & true\\
true & true & true & true\\
\end{tabularx}
\vspace{3mm}\footnotesize These distinctions determine which operation or control structure fits the problem.
\note{\par Syllabus reading: Deitel Ch 7 (as listed). CLO-2.}
\end{frame}

\begin{frame}[fragile,t]{First example: execution}
\begin{lstlisting}
int marks = 50;
boolean valid = marks >= 0 && marks <= 100;
boolean passed = valid && marks >= 50;
System.out.println(valid);
System.out.println(passed);
\end{lstlisting}
\note{This is the main body from Java\_Examples/Lecture07.java. The source file includes the complete class. Predict the result before the next slide.\par Syllabus reading: Deitel Ch 7 (as listed). CLO-2.}
\end{frame}

\begin{frame}[fragile,t]{First example: result}
\begin{columns}[T,onlytextwidth]\begin{column}{0.60\textwidth}
true\par true\par 50 satisfies the range and includes the passing boundary.
\end{column}\begin{column}{0.35\textwidth}\small
Intervals and precedence

Negation and boundary cases
\end{column}\end{columns}
\note{Expected output and interpretation: true
true
50 satisfies the range and includes the passing boundary.\par Syllabus reading: Deitel Ch 7 (as listed). CLO-2.}
\end{frame}

\begin{frame}[fragile,t]{Guided practice}
\begin{columns}[T,onlytextwidth]\begin{column}{0.60\textwidth}
A ticket discount applies to ages below 12 or at least 60. Write the condition and test ages 11, 12, 59 and 60.
\end{column}\begin{column}{0.35\textwidth}\small
Attempt the solution before continuing.

Use the definitions and examples from this lecture.
\end{column}\end{columns}
\note{Give students 8 minutes to attempt the problem. The following slides provide a complete solution. Original solution guidance: age \textless{} 12 || age \textgreater{}= 60
Results: true, false, false, true.\par Syllabus reading: Deitel Ch 7 (as listed). CLO-2.}
\end{frame}

\begin{frame}[fragile,t]{Solution strategy}
\begin{columns}[T,onlytextwidth]\begin{column}{0.60\textwidth}
\begin{itemize}
\item Express each qualifying age group as its own comparison.
\item Join the comparisons with ||.
\item Evaluate ages immediately below and at each boundary.
\end{itemize}
\end{column}\begin{column}{0.35\textwidth}\small
The next program uses fixed inputs so its result can be reproduced.
\end{column}\end{columns}
\note{Discuss why each step is needed. State the input assumptions before executing the program.\par Syllabus reading: Deitel Ch 7 (as listed). CLO-2.}
\end{frame}

\begin{frame}[fragile,t]{Complete solution}
\begin{lstlisting}
public class Solution07 {
  public static void main(String[] args) throws Exception {
    int age = 11;
    System.out.println(age < 12 || age >= 60);
    age = 12;
    System.out.println(age < 12 || age >= 60);
    age = 59;
    System.out.println(age < 12 || age >= 60);
    age = 60;
    System.out.println(age < 12 || age >= 60);
  }
}
\end{lstlisting}
\note{Complete runnable file: Java\_Examples/Solution07.java. Inputs are fixed in the program.  \par Syllabus reading: Deitel Ch 7 (as listed). CLO-2.}
\end{frame}

\begin{frame}[fragile,t]{Execution trace}
\small\renewcommand{\arraystretch}{1.5}
\rowcolors{2}{pale}{white}
\begin{tabularx}{\textwidth}{>{\raggedright\arraybackslash}X>{\raggedright\arraybackslash}X>{\raggedright\arraybackslash}X>{\raggedright\arraybackslash}X}
\rowcolor{navy}
\color{white}\bfseries Age & \color{white}\bfseries age \textless{} 12 & \color{white}\bfseries age \textgreater{}= 60 & \color{white}\bfseries Discount\\
11 & true & false & true\\
12 & false & false & false\\
59 & false & false & false\\
60 & false & true & true\\
\end{tabularx}
\vspace{3mm}\footnotesize Follow the changing state in execution order.
\note{Manually derived trace for the complete solution. Express each qualifying age group as its own comparison. Join the comparisons with ||. Evaluate ages immediately below and at each boundary.\par Syllabus reading: Deitel Ch 7 (as listed). CLO-2.}
\end{frame}

\begin{frame}[fragile,t]{Solution result}
\begin{columns}[T,onlytextwidth]\begin{column}{0.60\textwidth}
true\par false\par false\par true
\end{column}\begin{column}{0.35\textwidth}\small
Why this result follows

Evaluate ages immediately below and at each boundary.
\end{column}\end{columns}
\note{Expected result: true
false
false
true\par Syllabus reading: Deitel Ch 7 (as listed). CLO-2.}
\end{frame}

\begin{frame}[fragile,t]{A common incorrect approach}
boolean inRange = 0 \textless{}= marks \textless{}= 100;
\vspace{1mm}\begin{block}{Example}
\begin{lstlisting}
Use marks >= 0 && marks <= 100. Each comparison must produce a Boolean operand for &&.
\end{lstlisting}
\end{block}
\note{Ask students to predict the failure before discussing the explanation. The right side gives the correction.\par Syllabus reading: Deitel Ch 7 (as listed). CLO-2.}
\end{frame}

\begin{frame}[fragile,t]{Boundary and failure checks}
\begin{columns}[T,onlytextwidth]\begin{column}{0.60\textwidth}
Check the stated input assumptions.

Write a library eligibility condition from two age and membership rules of your choice. Supply a truth table.
\end{column}\begin{column}{0.35\textwidth}\small
age \textless{} 12 || age \textgreater{}= 60\par Results: true, false, false, true.
\end{column}\end{columns}
\note{Use the specification to derive each expected result. Exercise solution: age \textless{} 12 || age \textgreater{}= 60
Results: true, false, false, true.\par Syllabus reading: Deitel Ch 7 (as listed). CLO-2.}
\end{frame}

\begin{frame}[fragile,t]{Check your understanding}
\begin{columns}[T,onlytextwidth]\begin{column}{0.60\textwidth}
Evaluate true || false \&\& false. What is the opposite of score \textgreater{}= 50?
\end{column}\begin{column}{0.35\textwidth}\small
Write a prediction and explain the rule that supports it.
\end{column}\end{columns}
\note{Answer: true, because \&\& binds first. The opposite is score \textless{} 50.\par Syllabus reading: Deitel Ch 7 (as listed). CLO-2.}
\end{frame}

\begin{frame}[fragile,t]{Answer and explanation}
\begin{columns}[T,onlytextwidth]\begin{column}{0.60\textwidth}
true, because \&\& binds first. The opposite is score \textless{} 50.
\end{column}\begin{column}{0.35\textwidth}\small
Use marks \textgreater{}= 0 \&\& marks \textless{}= 100. Each comparison must produce a Boolean operand for \&\&.
\end{column}\end{columns}
\note{Reconnect the answer to the relevant definition rather than treating it as a fact to memorise.\par Syllabus reading: Deitel Ch 7 (as listed). CLO-2.}
\end{frame}


\begin{frame}[fragile,t]{Compare cases and predict outcomes}
\small\renewcommand{\arraystretch}{1.5}\rowcolors{2}{pale}{white}
\begin{tabularx}{\textwidth}{>{\raggedright\arraybackslash}X>{\raggedright\arraybackslash}X>{\raggedright\arraybackslash}X}\rowcolor{navy}
\color{white}\bfseries Age & \color{white}\bfseries Has card & \color{white}\bfseries age \textgreater{}= 18 \&\& hasCard\\
17 & true & false\\
18 & false & false\\
18 & true & true\\
\end{tabularx}\par\vspace{3mm}\begin{exampleblock}{Discuss}Explain each expected result before running the example. Which case would reveal a wrong assumption?\end{exampleblock}
\note{Read the first two columns and invite predictions before discussing the outcome.}
\end{frame}

\begin{frame}[fragile,t]{Apply the idea}
\begin{alertblock}{Independent practice}Admission requires age at least 18 and a valid card. Write both an admitted condition and its logical negation. Test the age boundary.\end{alertblock}\vspace{3mm}\begin{enumerate}\item Work through the input by hand.\item Write or correct the program.\item Justify the expected result.\end{enumerate}\note{Allow 3–5 minutes, then compare answers in pairs. Admission requires age at least 18 and a valid card. Write both an admitted condition and its logical negation. Test the age boundary.}
\end{frame}

\begin{frame}[fragile,t]{Solution and reasoning}
\begin{lstlisting}
boolean admitted = age >= 18 && hasCard;
boolean denied = age < 18 || !hasCard;
\end{lstlisting}\vspace{1mm}\small\begin{itemize}
\item Negating AND changes it to OR and negates each operand.
\item At age 18, the card determines the result.
\item At age 17, admission is false even with a valid card.
\end{itemize}\note{Answer: Negating AND changes it to OR and negates each operand. At age 18, the card determines the result. At age 17, admission is false even with a valid card.}
\end{frame}

\begin{frame}[fragile,t]{Divisibility with AND, OR and XOR}
\begin{itemize}
\item Remainder zero expresses divisibility by a nonzero integer.
\item AND requires both properties; OR accepts either or both.
\item Boolean XOR is true when exactly one property is true.
\end{itemize}
\vspace{3mm}\footnotesize\textcolor{accent}{Source: supplied ISB campus slides. See speaker notes and ISB\_Source\_Map.md.}
\note{Adapted from the supplied ISB campus folder: Lecture{-}8{-}9.pptx, slides 31, 32, 33, 34, 35. Adapted explanation and runnable implementation; sample inputs and traces added for this course.}
\end{frame}

\begin{frame}[fragile,t]{Worked problem: Divisibility with AND, OR and XOR}
\begin{block}{Task}For 6, evaluate divisibility by 2 and 3 using all three Boolean combinations.\end{block}
\small Input: Values are fixed in the program for a reproducible demonstration.\par\vspace{3mm}\begin{exampleblock}{Before running}Predict the output and identify the variables that change.\end{exampleblock}
\note{Adapted from the supplied ISB campus folder: Lecture{-}8{-}9.pptx, slides 31, 32, 33, 34, 35. Adapted explanation and runnable implementation; sample inputs and traces added for this course.}
\end{frame}

\begin{frame}[fragile,t]{Complete ISB example (1/1)}
\begin{lstlisting}
public class IsbLecture07 {
  public static void main(String[] args) throws Exception {
    int number = 6;
    boolean by2 = number % 2 == 0;
    boolean by3 = number % 3 == 0;
    System.out.println(by2 && by3);
    System.out.println(by2 || by3);
    System.out.println(by2 ^ by3);
  }

}
\end{lstlisting}
\note{Adapted from the supplied ISB campus folder: Lecture{-}8{-}9.pptx, slides 31, 32, 33, 34, 35. Adapted explanation and runnable implementation; sample inputs and traces added for this course. Complete file: ISB\_Java\_Examples/IsbLecture07.java.}
\end{frame}

\begin{frame}[fragile,t]{Worked trace and expected result}
\small\renewcommand{\arraystretch}{1.4}\rowcolors{2}{pale}{white}
\begin{tabularx}{\textwidth}{>{\raggedright\arraybackslash}X>{\raggedright\arraybackslash}X>{\raggedright\arraybackslash}X}\rowcolor{navy}
\color{white}\bfseries Number & \color{white}\bfseries AND / OR & \color{white}\bfseries XOR\\
6 & true / true & false\\
4 & false / true & true\\
5 & false / false & false\\
\end{tabularx}\par\vspace{2mm}
\begin{block}{Expected console output}\ttfamily\footnotesize true\par true\par false\end{block}
\note{Adapted from the supplied ISB campus folder: Lecture{-}8{-}9.pptx, slides 31, 32, 33, 34, 35. Adapted explanation and runnable implementation; sample inputs and traces added for this course. Trace and expected values independently worked for this edition.}
\end{frame}

\begin{frame}[fragile,t]{Extend the ISB example}
\begin{alertblock}{Practice}Why would independent if statements be appropriate when displaying all properties that hold?\end{alertblock}\vspace{3mm}\begin{exampleblock}{Explain your reasoning}More than one property can hold. An if{-}else chain would suppress later matching properties.\end{exampleblock}
\note{Adapted from the supplied ISB campus folder: Lecture{-}8{-}9.pptx, slides 31, 32, 33, 34, 35. Adapted explanation and runnable implementation; sample inputs and traces added for this course. Answer: More than one property can hold. An if{-}else chain would suppress later matching properties.}
\end{frame}
\begin{frame}[fragile,t]{Lecture summary}
\begin{columns}[T,onlytextwidth]\begin{column}{0.60\textwidth}
\begin{itemize}
\item relational and logical conditions
\item Boolean expressions
\item explicit grouping for compound rules
\end{itemize}
\end{column}\begin{column}{0.35\textwidth}\small
\begin{itemize}
\item Express each qualifying age group as its own comparison.
\item Join the comparisons with ||.
\item Evaluate ages immediately below and at each boundary.
\end{itemize}
\end{column}\end{columns}
\note{Ask students to give one concrete example for the concepts listed. Use the worked solution to connect the topics.\par Syllabus reading: Deitel Ch 7 (as listed). CLO-2.}
\end{frame}

\begin{frame}[fragile,t]{After-class practice}
\begin{columns}[T,onlytextwidth]\begin{column}{0.60\textwidth}
Write a library eligibility condition from two age and membership rules of your choice. Supply a truth table.
\end{column}\begin{column}{0.35\textwidth}\small
Syllabus reading\par Deitel Ch 7 (as listed)

Next lecture: Selection with if and else
\end{column}\end{columns}
\note{Reading labels follow the supplied outline. Check topic headings against the available textbook edition. Require a program and expected results for the practice task.\par Syllabus reading: Deitel Ch 7 (as listed). CLO-2.}
\end{frame}
\end{document}
